% =====================================================================
\subsection{Multi-domain provisioning workflow}
\label{sec:workflow}
% =====================================================================

Clause 6 ties the information to the service life cycle. When the QKD
network is introduced, the orchestrator acquires its topology (§5.1);
during operation it receives status and inventory, on demand or
periodically (§5.2); after a service is activated it monitors its status
(§5.3); at any time it receives notifications (§5.5).

§6.2 describes the complete case (Fig.~\ref{fig:seq018}): an orchestrator
coordinates the QKD controller and the OTN controller, which do not talk to
each other directly.

% Figure: end-to-end QKD + OTN provisioning (adapted from GS QKD 018, Fig. 2)
\begin{figure}[t]
\centering
\begin{tikzpicture}[font=\scriptsize]
  \def\xA{0}\def\xB{4.0}\def\xC{8.4}\def\xD{12.2}
  \node[actor] at (\xA,0) {Administrator};
  \node[actor] at (\xB,0) {Orchestrator};
  \node[actor,fill=green!8] at (\xC,0) {QKD controller};
  \node[actor,fill=orange!10] at (\xD,0) {OTN controller};
  \foreach \x in {\xA,\xB,\xC,\xD} \draw[gray] (\x,-0.4) -- (\x,-9.0);
  \draw[msg] (\xA,-1.0) -- node[above]{service request (SAE, QoS)} (\xB,-1.0);
  \draw[msg] (\xB,-1.5) -- ++(0.4,0) -- ++(0,-0.4) -- (\xB,-1.9);
  \node[right,align=left] at (\xB+0.45,-1.7) {address matching\\OTN node $\leftrightarrow$ QKD node};
  \draw[blue!60,rounded corners] (\xB-0.3,-2.55) rectangle (\xC+0.3,-3.85);
  \node[blue!60,anchor=south west,font=\scriptsize\itshape] at (\xB-0.3,-2.55) {repeatable loop};
  \draw[msg] (\xB,-3.0) -- node[above]{service link creation + constraints} (\xC,-3.0);
  \draw[ret] (\xC,-3.6) -- node[above]{\texttt{svc\_link\_id}, candidate paths} (\xB,-3.6);
  \draw[msg] (\xB,-4.6) -- node[above]{deploy chosen path (\texttt{path\_id})} (\xC,-4.6);
  \draw[msg] (\xC,-5.0) -- ++(0.4,0) -- ++(0,-0.4) -- (\xC,-5.4);
  \node[right,align=left] at (\xC+0.45,-5.2) {key relay along\\the path};
  \draw[ret] (\xC,-6.1) -- node[above]{keys ready} (\xB,-6.1);
  \draw[msg] (\xB,-6.9) -- node[above]{application provisioning (\texttt{qkd\_app})} (\xC,-6.9);
  \draw[ret] (\xC,-7.6) -- node[above]{\texttt{sdqkdn\_application\_deployed}} (\xB,-7.6);
  \draw[msg] (\xB,-8.6) -- node[above,fill=white,inner sep=1pt]{OTN nodes may fetch keys from the QKD nodes} (\xD,-8.6);
\end{tikzpicture}
\caption[End-to-end provisioning of a QKD service across a QKD network and an OTN]{End-to-end provisioning of a QKD service across a QKD network and
an OTN (simplified from~\cite[Fig.~2]{qkd018}). The interaction with the
OTN controller is out of the scope of the specification.}
\label{fig:seq018}
\end{figure}


The steps are:
\begin{enumerate}
  \item the administrator requests a service from the orchestrator,
        specifying the SAEs in the OTN nodes and their QoS;
  \item the orchestrator identifies the controllers involved and performs
        address matching; it already knows the addresses of the QKD nodes
        thanks to topology and inventory;
  \item it asks the QKD controller to create a service link and compute
        paths, with the desired constraints;
  \item the controller assigns an \texttt{svc\_link\_id} and returns the
        candidate paths, each with a priority; this step may be repeated
        with different constraints;
  \item the orchestrator picks a path and requests its deployment; the
        nodes generate and relay keys along the path until both ends are
        ready;
  \item once confirmed, the orchestrator requests provisioning of the
        application on the service link and tells the OTN controller that
        the OTN nodes may start fetching keys from the designated QKD
        nodes.
\end{enumerate}
Note the parallel with the virtual-link workflow of 015: there the
controller computed and configured the path on its own, here it
\emph{proposes} it and waits for the orchestrator's choice. With several
QKD domains the orchestrator repeats steps 3–5 for each domain along the
chosen inter-domain path, joining the segments through border nodes: this
is what Sim et al.\ do (§\ref{sec:sim}).

% =====================================================================
\subsection{Protocols and security}
\label{sec:protocolli}
% =====================================================================

\paragraph{Transport (§8).}
Like 015, the standard separates the data model (YANG~1.1) from the
transport protocol, which can be NETCONF~\cite{rfc6241} or
RESTCONF~\cite{rfc8040}. With RESTCONF every resource is a URI, reads use
\texttt{GET} and configuration uses \texttt{POST}, \texttt{PUT},
\texttt{PATCH} and \texttt{DELETE}, with XML or JSON encoding;
notifications are received by subscribing to the corresponding streams.
The standard notes that the stateless REST paradigm suits this kind of API.
As an example, reading the topology and invoking an RPC would look like:
\begin{lstlisting}
GET  /restconf/data/etsi-qkd-sdn-topology:qkd_network
POST /restconf/operations/etsi-qkd-sdn-connectivity:create-service-link-id-for-app
\end{lstlisting}

\paragraph{Security (§7).}
The controller has a central view of the network and passes a subset of it
to the orchestrator; protection measures for the interface (authentication,
channel encryption, e.g.\ TLS under RESTCONF) are left to implementations.
The fundamental principle, inherited from GS~QKD~015, is that
\textbf{QKD-derived keys never pass through either the controller or the
orchestrator}: the YANG modules contain only key rates and the KME address.
A compromised orchestrator can therefore cause service disruption (wrong
paths, exhausted resources, denial of service) but cannot read keys. The
standard recommends careful design, installation and personalisation
procedures to mitigate malicious requests and attacks on control channels.
